\documentclass[]{deedy-resume-openfont}
\begin{document}

\namesection{Yash}{Akhauri\hspace{5mm}\includegraphics[width=1.5cm, height=1.5cm]{CV_YashAkhauri.png}}{
% \href{https://quirkyai.wordpress.com}{Blog  |  } 
\href{https://github.com/akhauriyash}{GitHub  |  } 
\href{https://www.linkedin.com/in/akhauriyash/}{LinkedIn  |  } 
\href{mailto:akhauri.yash@gmail.com}{akhauri.yash@gmail.com} | +91 78915 12802 
}
% \intro{A Research Scientist specializing in Neural Network optimization with a keen interest in Mathematics and Physics.}
% \vspace{5mm}

\section{Education}
\runsubsection{BITS Pilani}
\descript{| B.E. in Electronics and Instrumentation}
\location{Aug 2016 - May 2020 | RJ, India\\
\custombold{Fluent: Python3, PyTorch, Mathematica \\  
 Familiar: C++, Java, CUDA, OpenMP, Tensorflow, Android Studio, LibGDX, Docker}}
\sectionsep
\section{Experience}

\runsubsection{Intel}
\descript{| Research Scientist \hfill \texttt{May 2020 - Present, Bangalore, India}}
\begin{tightemize}
\item Research Scientist at the Cloud Systems Research (CSR) Lab in the Systems and Software Research (SSR) Group.
\end{tightemize}
\sectionsep

\runsubsection{Xilinx Research}
\descript{| Visiting Scholar \hfill \texttt{Aug 2019 - May 2020,\quad Dublin, Ireland}}
\begin{tightemize}
\item Developed a library for co-design of neural network topologies and reconfigurable hardware that maps to an efficient FPGA implementation without the need for a custom accelerator architecture or a scheduler.
\item Targeted the Jet Substructure Classification task as part of CERN LMS L1 trigger experiments, used the library to deploy models with $10\times$ lower latency than FPGA4HEP designs.Demonstrated quantization library Brevitas to CERN \href{https://github.com/akhauriyash/FPGA4HEP-Brevitas/}{\texttt{[GitHub]}}.
\end{tightemize}

\sectionsep
\runsubsection{Uraniom}
\descript{| Research Intern \hfill \texttt{Jan 2019 - Jul 2019,\quad  France (Remote)}}
\begin{tightemize}
\item Implementing semantic segmentation models for face transfer across GIFs, progressive GANs for realistic UV map generation and explored effective weight-sharing strategies for neural networks under a research collaboration.
\end{tightemize}

\sectionsep
\runsubsection{Wolfram}
\descript{| Undergraduate Researcher \hfill \texttt{June 2018 - July 2018,\quad  \hspace{4pt} Massachusetts}}
\begin{tightemize}
\item Developed HadaNet MLPs in the Wolfram Language and worked on C OpenMP kernels for GEMM and Convolutions using the Hadamard Binarization algorithm.  \href{https://github.com/akhauriyash/XNOR-Intel-ISA}{\texttt{[OpenMP kernel]}} | \href{https://github.com/akhauriyash/XNOR-convolution}{\texttt{[CUDA kernel]}} | \href{https://docs.google.com/document/d/18uynX2yDSWm1BVCtG3Rd4CRb6xHiRxbvprUBTb4lvjY/edit?usp=sharing}{\texttt{[Whitepaper]}}
\end{tightemize}

\sectionsep

\section{Research}


\subsection{Papers}
\descript{LogicNets: Co-Designed Neural Networks and Circuits for Extreme-Throughput Applications \\\href{https://arxiv.org/abs/2004.03021}{\texttt{[IEEE \footnote[1]{To be published}]}} {\small Yash Akhauri*, Yaman Umuroglu*, Nicholas J. Fraser, Michaela Blott} \hfill \texttt{\emph{FPL'20} | Sweden | Sept 2020}} 
Paper accepted at \custombold{FCCM'20} as a poster presentation. Video presentation can be found \href{https://www.youtube.com/watch?v=jJRwyHD_UUI}{here.}\\
\sectionsep 

\descript{High-Throughput DNN Inference with LogicNets\\\href{https://www.fccm.org/home/program/}{\texttt{[IEEE \footnote[1]{To be published}]}} {\small Yaman Umuroglu, Yash Akhauri, Nicholas J. Fraser, Michaela Blott} \hfill\texttt{\emph{FCCM'20} | AR, USA | May 2020}}
\sectionsep 

\descript{HadaNets: Flexible Quantization Strategies for Neural Networks\\ \href{https://ieeexplore.ieee.org/document/9025370/}{\texttt{[IEEE]}} {\small Yash Akhauri} \hfill \texttt{\emph{CVPR'19 Workshop} | CA, USA | Jun 2019}} 
Paper accepted at \custombold{CVPR'19 UAVision workshop - Orals}. \\ 
Delivered a "Theatre Talk" and poster at the Intel Demo Booth at CVPR'19.\\
\sectionsep

\descript{Exposing Hardware Building Blocks to Machine Learning Frameworks\\ \href{https://arxiv.org/abs/2004.05898}{\texttt{[ArXiv -- Bachelor's Thesis]}} {\small Yash Akhauri} \hfill \texttt{Dec 2019}} 
\sectionsep
\vspace{1mm}

\clearpage

\subsection{Talks}
\descript{Wolfram Technology Conference\\ {\small Speaker} \hfill \texttt{Champaign, IL | Oct. 2018}}
Delivered a talk on my research on Hadamard Neural Networks. \\

\sectionsep
\descript{Intel AI Meetup\\ {\small Speaker} \hfill \texttt{Delhi, IN | Sept. 2018}}
\href{https://docs.google.com/presentation/d/1PuxUM6RmDI5Dq7XgFoVWNxv_TOQ3aux8GiQ9aWlpUjg/edit?usp=sharing}{\texttt{[PPTX]}} \href{https://software.intel.com/en-us/articles/developing-hadamard-neural-networks-on-the-intel-xeon-scalable-processor}{\texttt{[Article]}}
Spoke about my research on scaling AI using Intel technologies. This event was organized by Intel. \\

\sectionsep
\descript{Intel AI DevCon\\ {\small Poster} \hfill \texttt{San Francisco, Bangalore | May \& Aug 2018}}
Presented posters on quantized GEMM kernels for Intel Xeon Phi\\
\vspace{1mm}
\sectionsep


\subsection{Grants}

\descript{Intel Nervana Early Innovators Grant \hfill \texttt{\$5000}}
Received research grant to develop Binary Precision Neural Networks and Real time Artistic Style Transfer. The technical article can be found  \href{https://software.intel.com/en-us/articles/art-em-artistic-style-transfer-to-virtual-reality-final-update}{\texttt{[here.]}}
The code can be found \href{https://github.com/akhauriyash/Fast-style-transfer}{\texttt{[here.]}}\\
Intel AI Academy Success Story \href{https://software.intel.com/en-us/articles/developing-hadamard-neural-networks-on-the-intel-xeon-scalable-processor}{\texttt{[Link]}} published by Intel for the research done as part of this grant in the field of Quantized Neural Networks.
\sectionsep 

\descript{Intel CVPR Travel Grant \hfill \texttt{\$3000}}
Received a travel grant from Intel to present research at the Intel Demo Booth at CVPR'19.\\
\sectionsep

\descript{Wolfram Student Aid \hfill \texttt{\$2400}}
Received aid to attend the Wolfram Summer School and develop Hadamard Binary Neural Networks. \\
\sectionsep

\descript{KVPY Scholar}
Selected as a KVPY scholar by the Department of Science and Technology, Government of India. \\
\sectionsep

\descript{INSPIRE Scholarship}
Selected for the INSPIRE Scholarship by the Department of Science and Technology (DST), Government of India.\\
\vspace{1mm}
\sectionsep

\subsection{Projects}

\descript{Exploiting Huffman Coding and Weight-Sharing for Memory-Efficient Inference on FPGAs. \hfill \texttt{PyTorch}}
Developed a neural network weight sharing strategy and proposed a methodology to leverage this strategy in the FINN architecture for Multi-Layer Offload and Data-Flow Architectures of Neural Network deployment on FPGAs. \\
\sectionsep

\descript{Whitepaper - Improving distributed mesh computing with Hadamard Binary Neural Networks.}
\href{https://docs.google.com/document/d/18uynX2yDSWm1BVCtG3Rd4CRb6xHiRxbvprUBTb4lvjY/edit?usp=sharing}{\texttt{[Link]}}\\
\sectionsep

\descript{xGEMM \& xCONV \hfill \texttt{C++, CUDA, OpenMP}}
Coded efficient 3D convolutional and GEMM kernels for XNOR (bit quantized) networks using CUDA C programming and OpenMP. Optimized kernels are for Intel processors and Nvidia GPUs. Invited to present a poster at Intel AI DevCon and Intel AI Student Ambassador Summit, San Francisco. The codes can be found here: \href{https://github.com/akhauriyash/XNOR-Intel-ISA}{\texttt{[OpenMP kernel]}} | \href{https://github.com/akhauriyash/XNOR-convolution}{\texttt{[CUDA kernel]}}.\\
\sectionsep

% \descript{GEMM: From Pure C to SSE Optimized Micro Kernels \hfill \texttt{C++}}
% Followed 
%  \href{http://apfel.mathematik.uni-ulm.de/~lehn/sghpc/gemm/}{\texttt{this tutorial}} to build an implementation of GEMM that achieves performance close to that of the BLIS kernels. \\
% \sectionsep

\descript{Real time artistic style transfer \hfill \texttt{Python, Tensorflow, OpenCV}}
Worked on pruned generator networks with Adaptive Instance Normalization to give stylization speeds of 25FPS at VR resolution. Integrated pruned network with Android IP Webcam for live stylization of camera feed. The article can be found \href{https://software.intel.com/content/www/us/en/develop/articles/art-em-artistic-style-transfer-to-virtual-reality-final-update.html}{\texttt{[here]}}. 
\sectionsep

\descript{GravDash \hfill \texttt{Java, LibGDX, Android Studio}}
Developed an android game using Java and libGDX as the framework. \\
\sectionsep

\descript{Blog \hfill  \texttt{Python, Java, C++, Tensorflow, OpenMP}}
Maintaining a \href{https://quirkyai.wordpress.com}{\texttt{[blog]}} covering various AI   related topics with over 10000 hits. \\
\sectionsep


% \sectionsep
\end{document}
